\vspace{-0.5em}
\section{Introduction}

Artificial Intelligence (AI) systems, particularly Neural Networks (NNs) are being employed in critical sectors such as healthcare to interpret clinical images or in autonomous driving to recognize physical elements in traffic images. While highly performant, they often operate as black boxes that offer limited insight into the reliability of their outputs. This opacity becomes critical under adversarial, uncertain, or degraded input conditions. Conventional performance metrics such as accuracy or precision do not capture uncertainty or reliability, which can mislead decision-makers about the quality of a model’s outputs~\cite{10.1145/3290605.3300509}. For instance, a model may report 95\% accuracy in a classification task while being evaluated on mislabeled data. This accuracy evaluation does not account for dataset bias, label noise, or adversarial corruption. These challenges underscore the need for reliability measures that go beyond predictive confidence and that consider the quality and provenance of inputs, the reliability of training data, and the fit of learned parameters.

As AI systems become integral to critical applications, \emph{trustworthiness} has emerged as a fundamental requirement. A trustworthy system must be lawful, ethical, and robust~\cite{ec2019ethics}, some of which translate into measurable properties such as accuracy, robustness, fairness, and explainability~\cite{10.3389/fdata.2024.1467222}. Embedding these across the AI pipeline, from data collection to deployment, is essential for preventing failures caused by poor data quality, bias, or unstable training. 

Despite this growing awareness, most models still provide limited mechanisms for representing confidence in their outputs. Neural networks are typically trained with the assumption of clean, unambiguous data labels and thus lack in the capability to generalize across previously unseen borderline cases. As a result, predicted probabilities tend to reflect similarity to seen patterns rather than genuine uncertainty, and they often assume clean, in-distribution inference inputs. These assumptions are rarely satisfied in real-world or adversarial settings, where data quality, distributional shifts, and how well parameters are learned all influence prediction reliability. Consequently, existing approaches lack mechanisms to evaluate how uncertainty and trustworthiness in training inputs, activations, and parameters jointly affect model behavior. Even attribution-based methods (like SmoothGrad~\cite{smilkov2017smoothgradremovingnoiseadding}) fail to capture how training input, intermediate, and parameter trustworthiness jointly influence the output reliability. This gap motivates the development of tools that can explicitly reason about trustworthiness across the entire model lifecycle.

\paragraph*{Problem Statement}
Existing methods for trustworthiness and uncertainty estimation in neural networks face three key limitations. They \emph{neglect data provenance and quality}, assuming the training data are fully reliable; they offer only \emph{limited holistic propagation}, assessing reliability at the output layer alone; and they \emph{lack interpretability}, rarely yielding estimates that are both faithful to the model's reasoning and understandable to end users. Together these hinder reliable trustworthiness assessment in safety-critical or adversarial contexts.

\paragraph*{Contributions}
To address these challenges, we propose the \emph{Parallel Trust Assessment System (PaTAS)}, a framework for modeling and propagating trustworthiness in NNs using Subjective Logic (SL). Our contributions are:
\begin{enumerate}
    \item \emph{Parallel Trust Computation:} \textit{Trust Nodes} and \textit{Trust Functions} that mirror neural computations, enabling principled trust propagation during training and inference through SL discounting and fusion.
    \item \emph{Parameter Trust Update:} an algorithm that derives trust in learned parameters from gradient values, input trust, and label trust, aligning parameter reliability with the learning dynamics.
    \item \emph{Inference-Path Trust Assessment (IPTA):} a context-aware trust function using activation-path information to compute per-instance trust scores.
    \item \emph{Empirical Validation:} evaluation on real-world and adversarial datasets, demonstrating that PaTAS produces interpretable trust estimates that expose reliability gaps under poisoned, biased, and uncertain data.
\end{enumerate}

\paragraph*{Structure of the Paper}
The remainder of this paper is organized as follows: \cref{sec:back} introduces background on Subjective Logic and dataset trustworthiness assessment, \cref{sec:rw} reviews related work, \cref{sec:method} formalizes trust propagation in neural networks, \cref{sec:patas} details the PaTAS architecture, \cref{sec:evaluation} presents experiments, \cref{sec:discussion} discusses implications, and \cref{sec:conclusion} concludes the paper and provides an outlook on future work.
